<TODO type="unknown-block-type"> <!--todo: unknown block type "html"--> <pre>&lt;p&gt; </pre> </TODO>
    \documentclass{article}
</p>
\pagenumbering{roman}
\usepackage[utf8]{inputenc}
\usepackage{geometry}
 \geometry{
 letterpaper,
 % total={170mm,257mm},
 left=20mm,
 right=20mm,
 top=20mm,
 }
 \usepackage{graphicx}
 \usepackage{titling}
 \usepackage{amsfonts}
\usepackage{fancyhdr}
\usepackage{amsmath}
\usepackage{amssymb}
\usepackage{graphicx}
\usepackage{enumerate}
\usepackage{verbatim}
\usepackage{setspace}
\usepackage{pdflscape}
\usepackage{float}
%\usepackage[margin=1in, left = 1.5in]{geometry}
\usepackage{hyperref}
\usepackage{tikz}
\usetikzlibrary{matrix,arrows}
\usepackage{tikz-cd}
\usepackage{multicol}
\usepackage{color}
\usepackage{bbding}
\usepackage{cite}
\usepackage{mathtools}
\usepackage{pdfpages}

\usepackage[capitalise]{cleveref}


\newtheorem{thm}{Theorem}[section]
\newtheorem{lema}[thm]{Lemma}
\newtheorem{prop}[thm]{Proposition}
\newtheorem{cor}[thm]{Corollary}
\newtheorem{conjecture}[thm]{Conjecture}
\newtheorem{prob}{Problem}
%\newtheorem{question}{Question}[section]
\newtheorem{ex}[thm]{Example}
\newtheorem{rmk}[thm]{Remark}
\newtheorem{defn}[thm]{Definition}
\newtheorem{qu}[thm]{Question}
\newtheorem{ans}[thm]{Answer}
\newtheorem{note}[thm]{Note}
\newtheorem{hw}{Homework}

\newcommand{\qedsymbol}{\openbox}
%\renewcommand{\baselinestretch}{1.5}

\newcommand{\Mod}[1]{\ \left(\mathrm{mod}\ #1\right)}
\newcommand{\rk}{\textrm{rk}}
\newcommand{\limx}{\lim_{x\rightarrow\infty}}
\newcommand{\qed}{\hfill\qedsymbol}
\newcommand{\vsp}{\vspace{1em}}
\newcommand{\pf}{\textit{\underline{Pf:}} }
\newcommand{\soln}{\textit{\underline{Soln:}} }

\newcommand{\CC}{\mathbb{C}}
\newcommand{\PP}{\mathbb{P}}
\newcommand{\NN}{\mathbb{N}}
\newcommand{\QQ}{\mathbb{Q}}
\newcommand{\RR}{\mathbb{R}}
\newcommand{\ZZ}{\mathbb{Z}}
\newcommand{\FF}{\mathbb{F}}
\newcommand{\TT}{\mathbb{T}}
\newcommand{\M}{\overline{M}}
\newcommand{\fpn}{\mathbb{F}_{p^n}}
\newcommand{\fp}{\mathbb{F}_{p}}

\newcommand{\stuff}{{\color{white}nothing so there's a new line}}

\newcommand{\sn}{(s_n)}
\newcommand{\tn}{(t_n)}
\newcommand{\an}{(a_n)}
\newcommand{\bn}{(b_n)}
\newcommand{\cn}{(c_n)}

\newcommand{\suman}{\sum_{n=1}^\infty a_n}
\newcommand{\sumbn}{\sum_{n=1}^\infty b_n}
\newcommand{\psa}{\sum_{n=0}^\infty a_n x^n}
\newcommand{\nosum}{\sum_{n=0}^\infty }
\newcommand{\onesum}{\sum_{n=1}^\infty }
\newcommand{\inv}{^{-1}}
\newcommand{\vecv}[1]{\vec{v_{#1}}}
\newcommand{\<}{\langle}
\renewcommand{\>}{\rangle}
\newcommand{\ul}[1]{\underline{\hspace{#1in}}}
\newcommand{\lmx}[1]{\displaystyle\lim_{x\rightarrow #1}}

 \title{Worksheet 1: Sections 1.1, 1.2, 2.1 - Review of Functions, Rates of Change, Secant and Tangent
Lines
}
\author{MATH 1551 - Worksheets}
\date{Fall 2024}
 
 \usepackage{fancyhdr}
\fancypagestyle{plain}{%  the preset of fancyhdr 
    \fancyhf{} % clear all header and footer fields

    % \fancyfoot[L]{\thedate}
    \fancyhead[L]{\theauthor}
    \fancyhead[R]{Worksheet 1}
  
}
\pagestyle{plain}
\makeatletter
\def\@maketitle{%
  \newpage
  \null
  \vskip 1em%
  \begin{center}%
  \let \footnote \thanks
    {\LARGE \@title \par}%
    \vskip 1em%
    %{\large \@date}%
  \end{center}%
  \par
  \vskip 1em}
\makeatother

\usepackage{lipsum}  
\usepackage{cmbright}
%\pagestyle{empty}
\pagestyle{fancy}
\begin{document}

\maketitle

For reference as you go:



\begin{itemize}

\item [a)] The worksheet has two pages and includes six problems, some of which, labeled with (\textbf{AP}), are intended as additional practice outside of class.\medskip
\item[b)] The point-slope equation of a line with slope $m$ passing through the point $(x_1,y_1)$ is given by 
$$y-y_1 = m(x-x_1).$$
\iffalse 
\item Some useful identities:
    \begin{align*}
        a^2-b^2&=(a-b)(a+b) \\ \\
        (a-b)(a+b)&=a^2-b^2 \\ \\
        (a+b)^2&=a^2+2ab+b^2 \\ \\
        a^2+2ab+b^2 &= (a+b)^2 \\ \\
         (a-b)^2&=a^2-2ab+b^2 \\ \\
        a^2-2ab+b^2 &= (a-b)^2 \\ \\
        a(b+c)&=ab+bc \\ \\
        ab+bc&=a(b+c) \\ \\
    \end{align*}
    \item Factoring quadratics:
    \begin{align*}
    ax^2+bx+c = a(x-x_1)(x-x_2),
\end{align*}
where 
\begin{align*}
    x_1 = \frac{-b+\sqrt{b^2-4ac}}{2a}, \\ \\
    x_2 = \frac{-b-\sqrt{b^2-4ac}}{2a}. 
\end{align*}
\fi 
\item[c)] Completing the square:
\begin{align*}
    ax^2+bx+c = a(x+k)^2+h,
\end{align*}
    \begin{align*}
 \mbox{ where } \qquad   k= \frac{b}{2a} \qquad \mbox { and } \qquad 
    h=c-\frac{b^2}{4a}.
\end{align*}
    
\item [d)] Absolute value function:
$$f(x)=|x|=\begin{cases}
x\quad \textrm{ if }x \geq 0 \\ \\
-x\quad \textrm{ if } x < 0
\end{cases}$$

\item [e)] Some algebraic identities: 
\begin{align*}
    a^2-b^2&=(a-b)(a+b) \\ \\
    (a-b)(a+b) &=a^2-b^2
\end{align*}
\end{itemize}

\noindent\rule{\textwidth}{.5pt}
%\vfill
\begin{enumerate}
\iffalse 
\item Simplify each expression as far as possible, or state that the 
expression cannot be simplified.

\begin{enumerate}
    \item $\displaystyle {{x}\over{x+2}}$ \medskip
    \item $\displaystyle {{x^3-4x}\over{x+2}}$ \medskip
    \item $\displaystyle{{x^2+2x}\over{x^3-x^2-6x}}$ \medskip
    \item $\sqrt{x^2-2x+1}$  \medskip
    \item $\sqrt{x^2-1}$ \medskip
    \item $\sqrt{a^2 + b^2}$ \medskip
    \item $\sqrt{x^2}$ \medskip
\end{enumerate}

\item Let $f(x) = 3x^2 - 2x + 1$.  Find $f(2+h)$. \medskip

\fi
%\item Find the domain of the function $\displaystyle f(x) = \sqrt{ {{x+2}\over{x-3}} }$. \medskip

%\item Find the range of the function $y=x^2 - 6x + 5$ \medskip



%\item Let $\displaystyle f(x)=1+{{1}\over{x}}$ and $\displaystyle g(x)={{x}\over{2+x}}$.  \medskip

%Find the functions $f+g$, $fg$, $\displaystyle {{f}\over{g}}$, $g \circ f$, $f \circ g$, $g \circ g$, $f \circ f$ and determine the domain of each function.

   \item Find the domain and range of each of the following functions:
    \begin{multicols}{2}
    \begin{enumerate}
        \item (\textbf{AP}) $f(x)=x^3+3$  \medskip
        \item $g(x)=x^2 - 6x + 5$ \medskip
        \item (\textbf{AP}) $h(x)=2\sqrt{x+1}+1$ \medskip
        \item $k(x)=|x-2|-3$ \medskip
        \item (\textbf{AP}) $l(x)=-2\cos(x-\pi)$ \medskip
        \item $\displaystyle m(x)=\frac{\sin^2(2x)}{2}$ \medskip
        \item $\displaystyle y=\frac{2-x}{|3x-6|}$ \medskip
        \item (\textbf{AP}) $\displaystyle y=\frac{2}{\sqrt[4]{t}}$ \medskip
        \item (\textbf{AP}) $y=\ln(s-5)$ \medskip
        \item $y= -e^s+2$ 
    \end{enumerate}
    \end{multicols}\vsp 

\item Suppose that $f(x)=\sqrt{2x-1}$ and $g(x)=x^2+1$. Find $(f\circ g)(x)$, $(g \circ f)(x)$, and their domains. \vsp 
    \item For each of the following functions and intervals, find the difference quotient $$\frac{y_2-y_1}{x_2-x_1}.$$
    \begin{enumerate}
        \item $f(x) = \ln(x)+5$ on $[1,9]$. Simplify your answer, if possible. \medskip
        \item \textbf(\textbf{AP}) $g(x) = 4x^2-e^x$ on $[x,a]$, where $a > x$. \qquad \textit{(Note: No need to simplify your answer)} \medskip
        \item $k(x) = \sqrt{x-4}$ on $[x,x+h]$, where $h > 0$. \qquad \textit{(Note: No need to simplify your answer)} \medskip
    \end{enumerate}\vsp 

%\newpage 
    \item You may use Desmos to recall the graph of $\displaystyle f(x)=\frac{1}{x}$. Using what you know about reflecting, shifting, and stretching, match the transformations of $f(x)$ shown below with their equations. 
    \begin{align*}
         (a) \quad \displaystyle y=\frac{-3}{x}\qquad \qquad \qquad \qquad
        (b) \quad \displaystyle y=\frac{1}{x+1}-1 \qquad \qquad \qquad  \qquad
        (c) \quad \displaystyle y=2\left(\frac{1}{x+1}+1\right)
    \end{align*}
    \iffalse 
    \begin{enumerate}
        \item $\displaystyle y=\frac{-3}{x}$ \medskip
        \item $\displaystyle y=\frac{1}{x+1}-1$ \medskip
        \item $\displaystyle y=2\left(\frac{1}{x+1}+1\right)$ \medskip
    \end{enumerate}
    \fi
    \begin{figure}[h]
        \centering
        \includegraphics[scale=0.4]{WS1.PNG}
    \end{figure}\vsp 

\item  The graph of a function $f(x)$ %f(x)=-x^3+3x$ 
is shown in the figure. By inspecting the graph, fill in the bubble next to the correct ordering of the quantities $A$ through $E$. 

 \begin{minipage}[t]{0.6\textwidth}
        \vspace{0pt}
         \begin{itemize}
         \item [A.] The slope of the tangent line to the graph of $f(x)$ at $x=1$ \medskip
      \item [B.] The slope of the secant line between $x=0$ and $x=1$ \medskip
      \item [C.] The instantaneous rate of change of $f(x)$ at $x=-2$ \medskip
      \item [D.] The average rate of change of $f(x)$ over the interval $[-1,2]$ 
\medskip
\item [E.] $\displaystyle \frac{f(2)-f(-2)}{4}$
    \end{itemize}
    \item[] \textit{Answer:} 
        
    \end{minipage}%
    %\hfill
    \begin{minipage}[t]{0.4\textwidth}
\vspace{0pt}
        
      
        \includegraphics[scale=.2]{fa24quiz1pic1.png}
 
    \end{minipage} 
    \vsp 
    \item Find the equation of the secant line through the points on $g(x)=e^{x-1}-x-1$ where $x=1$ and $x=2$. \vsp 
\iffalse 
    
    \item 
    The graph of a function $f(x)=-x^3+3x$ is shown in the figure.
  \begin{center}
       \includegraphics[scale=0.2]{fa24quiz1pic1.png}
  \end{center}
  
    By inspecting the graph, fill in the bubble next to the correct ordering of the quantities $A$ through $D$. 
  \begin{itemize}
      \item [A.] The slope of the tangent line to the graph of $f(x)$ at $x=1$ \medskip
      \item [B.] The slope of the secant line between $x=0$ and $x=1$ \medskip
      \item [C.] The instantaneous rate of change of $f(x)$ at $x=-2$ \medskip
      \item [D.] The average rate of change of $f(x)$ over the interval $[-1,2]$ \medskip
  \end{itemize}
  \fi 
  \iffalse 
  \noindent\rule{\textwidth}{.5pt}
  
\textit{Extra Problems:} 
\begin{enumerate}
   
\item Use Desmos to graph the piecewise function: $f(x)=\begin{cases}
{\displaystyle \frac{1}{x+2}}\quad \textrm{ if }x > 0 \\ \\
|2x+1| \quad \textrm{ if } -2 \le x \le 0
\\ \\
\sqrt{5-x} \quad \textrm{ if } x < -2
\end{cases}$

Find the domain and range of $f(x)$. \medskip

\item Suppose we are interested in the instantaneous velocity of a rocket whose height is modeled by the function $y=98-2(t+4)^2$ at the instant where $t=2$ seconds. Find the average velocity of the rocket over a reasonable interval in order to approximate the instantaneous velocity. How accurate do you think your approximation is? Why? 

\end{enumerate}

\fi 
\iffalse 
%\item For each of the following functions, identify any asymptotes, the intervals on which the function is increasing/decreasing, and intervals of concavity, along with anything else you think may be helpful (like $x-$ or $y-$ intercepts). Use that information to sketch a graph of the function. 
   % \begin{enumerate}
      %  \item $f(x)=\frac{e^x}{x}$\vsp 
     %   \item $g(x)=x^5-x$\vsp 
     %   \item $h(x)=\tan^2(x)$
   % \end{enumerate}\vsp 


%\item Consider the Mean Value Theorem (MVT) (writing it down carefully now will help you in the next few problems too).
%\begin{enumerate}
   % \item Given a function $f(x)$ and a closed interval $[a,b],$ two conditions must be met in order for the MVT to apply. \medskip
   % \begin{enumerate}
      %  \item What is the first condition? \medskip
      %  \item What is the second condition? \medskip
   % \end{enumerate}

   % \item Now  consider Rolle's Theorem.
    %\begin{enumerate}
      %  \item Rolle's Theorem requires a third condition to be met in addition to the conditions of the MVT. What is it? \medskip
      %  \item Explain how Rolle's Theorem is related to the MVT.
   % \end{enumerate}
%\end{enumerate}\vsp 

\item Consider the following scenario, and decide whether you should use Rolle's Theorem or  the MVT to solve the problem, or whether neither applies. 

        \emph{You ride a snowboard along a hill shaped like a cosine wave, as pictured below. If you ride from $x=0$ to $x=2\pi$, argue that there must be a point where your snowboard is perfectly level.}
        \begin{figure}[h]
            \centering
\includegraphics[scale=0.25]{WS10pic1.png}
        \end{figure}


\item Decide whether the MVT applies to each of the following functions $f(x)$ over the specified interval $I$. Be sure to justify your answer! 

If the MVT does apply, find the point(s) guaranteed by the theorem, that is, the point(s) $c$ where $$f'(c)=\frac{f(b)-f(a)}{b-a}.$$

\begin{multicols}{2}
\begin{enumerate}
    \item $\displaystyle f(x)=\frac{1}{x}$, $I=[-1,1]$\medskip
    
    \item $\displaystyle  f(x)=\frac{1}{x}$, $I=[1,2]$\medskip

    \item $\displaystyle  f(x)=2\,e^x-2x^{1/3}$, $I=[-1,1]$ \medskip
    
    \item $\displaystyle  f(x)=|x|$, $I=[-1,1]$ \medskip
    
    \item $\displaystyle f(x)=|x|$, $I=[1,2]$\medskip

    \item $\displaystyle  f(x)=\sin\left(x/2\right)$, $I=[0,\pi]$ \medskip
\end{enumerate}
\end{multicols}
\vsp 

\item Decide whether each of the following statements is true or false. 

\emph{Remember: a logical statement is only true if it is ALWAYS true; if it is ever false, the implication is false.}

\begin{enumerate}
    \item If $f(x)$ is NOT continuous over $[a,b],$ then there is no value $x=c$ where $\displaystyle f'(c)=\frac{f(b)-f(a)}{b-a}.$ \vsp

    \item If $f(x)$ is differentiable over $[a,b],$ then there is at least one value $x=c$ where $\displaystyle f'(c)=\frac{f(b)-f(a)}{b-a}.$ \vsp
    
    \item If $f(x)$ is continuous on $[a,b]$ and differentiable on $(a,b),$ then there is exactly one  value $x=c$ where $\displaystyle f'(c)=\frac{f(b)-f(a)}{b-a}.$ \vsp
    
    \item If $f(x)$ is continuous over $[a,b],$ $f(x)$ takes on every value between $a$ and $b$ as $x$ moves from $a$ to $b$. \vsp
    
    \item If $f(x)$ is continuous over $[a,b],$ $f(x)$ takes on every value between $f(a)$ and $f(b)$ as $x$ moves from $a$ to $b$. \vsp
    
    \item If $f(x)$ is continuous over $[a,b]$ but not differentiable over $(a,b),$ then $f(x)$ may not take on every value between
$f(a)$ and $f(b)$ as $x$ moves from $a$ to $b$.
\end{enumerate}
\vsp



\item Compute the following antiderivatives.

\begin{multicols}{3}
    \begin{enumerate}
        \item $\displaystyle \int x(x^2+5) \,dx$\vsp 
        
        \item $\displaystyle \int \left(\frac{u}{\sqrt[3]{u}} - \frac{2}{u} \right)\,du$\vsp 
        
        \item $\displaystyle \int (3e^{5u}-u^e) \,du$\vsp 
        
      %  \item $\displaystyle \int \frac{\pi}{7x} \,dx$\vsp 
        
       % \item $\displaystyle \int (3e^\pi+ey+\pi^y) dy$\vsp 
        
        \item $\displaystyle \int \left( \sin(2x)+\frac{3}{\sqrt{1-4x^2}} \right) \,dx$\vsp 
        
        \item $\displaystyle \int (4^x\cdot 3^x+\pi) \,dx$\vsp 
        
        \item $\displaystyle \int \csc(3x)\cot(3x) \,dx$
    \end{enumerate}
\end{multicols} 


\fi
\end{enumerate}
\end{document}